% 分野別類題: 連立線形微分方程式（dx/dt = Ax）
% 係数行列の固有値・固有ベクトルによる解法。実固有値・複素固有値・初期値問題を扱う 3 問。
% コンパイル: TEXINPUTS="../../../00_common:$TEXINPUTS" latexmk -xelatex problems.tex
\documentclass[a4paper,11pt]{article}
\usepackage{preamble}

\begin{document}

\kamokumei{類題演習 --- 連立線形微分方程式}

\shijibun{次の連立常微分方程式を解き、導出過程を示せ。係数行列 $A$ の固有値・固有ベクトルを求め、それらを用いて一般解を構成すること。}

\shomon{問1.}{次の連立微分方程式の一般解を $x(t)$, $y(t)$ の形で求めよ。}
\[
\frac{dx}{dt} = x + 2y, \qquad \frac{dy}{dt} = 2x + y
\]

\shomon{問2.}{次の連立微分方程式の一般解を求めよ。}
\[
\frac{dx}{dt} = x - y, \qquad \frac{dy}{dt} = x + y
\]

\shomon{問3.}{次の初期値問題を解け。}
\[
\frac{dx}{dt} = 4x - 2y, \qquad \frac{dy}{dt} = x + y, \qquad x(0) = 1,\ \ y(0) = 0
\]

\end{document}
